% Part: first-order-logic
% Chapter: axiomatic-deduction
% Section: provability-quantifiers

% verification of properties of provability needed for maximally
% consistent sets in the completeness chapter.

\documentclass[../../../include/open-logic-section]{subfiles}

\begin{document}

\olfileid{fol}{axd}{qpr}

\olsection{\usetoken{S}{derivability} en de kwantoren (‘voor alle’ en ‘er bestaat’)}

\begin{explain}
  Voor het bewijs van de volledigheidsstelling hebben we ook de feiten
  over~$\Proves$ nodig die we hier met axiomatische afleidingen bewijzen.
\end{explain}

\begin{thm}
\ollabel{thm:strong-generalization} Stel dat $c$ een !!a{constant} is die
niet voorkomt in $\Gamma$ en ook niet in $!A(x)$. Als
$\Gamma \Proves !A(c)$, dan geldt
$\Gamma \Proves \lforall[x][!A(x)]$.
\end{thm}

\begin{proof}
Gebruik eerst de deductiestelling. Dan krijgen we
$\Gamma \Proves \ltrue \lif !A(c)$. De constante $c$ komt niet voor
in $\Gamma$ en ook niet in~$\top$. Daarom geldt
$\Gamma \Proves \ltrue \lif \lforall[x][!A(x)]$. Pas de
deductiestelling nogmaals toe. Dan krijgen we
$\Gamma \Proves \lforall[x][!A(x)]$.
\end{proof}

\begin{prop}
\ollabel{prop:provability-quantifiers}
\begin{tagenumerate}{prvEx,prvAll}
\tagitem{prvEx}{$!A(t) \Proves \lexists[x][!A(x)]$.}{}

\tagitem{prvAll}{$\lforall[x][!A(x)] \Proves !A(t)$.}{}
\end{tagenumerate}
\end{prop}

\begin{proof}
\begin{tagenumerate}{prvEx,prvAll}
\tagitem{prvEx}{Gebruik \olref[qua]{ax:q2} en daarna de deductiestelling.}{}
\tagitem{prvAll}{Gebruik \olref[qua]{ax:q1} en daarna de deductiestelling.}{}
\end{tagenumerate}
\end{proof}

\end{document}
